\section{The Contract and Its Execution Boundary}\label{sec:r37-institution}
\subsection{A saturated renewal agreement}
At the renewal date, a public branch $j\in\{1,\ldots,k\}$ is observed with probability $\pi_j>0$, where $\sum_j\pi_j=1$. Its continuation cap is $0<b_1<\cdots<b_k<1$. The provider chooses an intermediate tier $Y_j\in[0,1]$ and a terminal tier $C_j\in[0,1]$. All branches subsequently reach the same terminal public vertex. Unit payment coefficients and unit discounting make total remaining payment $Y_j+C_j$. A fixed service-benefit continuation minus the customer's outside continuation gives the cap. At the root the provider has committed to
\begin{equation}\label{eq:r37-saturation}
 \sum_j\pi_j\E[Y_j+C_j]=\bar b,
 \qquad \bar b:=\sum_j\pi_jb_j.
\end{equation}
Under expected participation, $\E[Y_j+C_j]\leq b_j$ for every $j$, with expectation conditional on that observed branch and taken before the internal draw. Positivity of the probabilities and \eqref{eq:r37-saturation} imply equality separately on every branch. Under pathwise participation, the stronger constraint is $Y_j+C_j\leq b_j$ almost surely. The same argument then implies equality almost surely, not merely in expectation.

The operating payoff is $f(C_j)-h_j(Y_j)$, with increasing concave $f$ and convex nondecreasing $h_j$ satisfying $h_j(0)=0$. General shape assumptions suffice for the deterministic and pathwise results. The global continuous algorithm specializes to the quadratic primitives in \eqref{eq:r34-primitives}. The negative intermediate contribution represents the net cost of an advance or expedited tier, while the terminal tier contributes net operating return. It does not mean that the customer has negative service benefit: that benefit has already entered the continuation cap. There are no endogenous public transitions, coupled commitments, or switching costs in the renewal theorem.

\subsection{What a writable code means}
At the observed branch, an encoder selects a symbol $Z\in\{1,\ldots,m\}$. At the terminal vertex the decoder observes only the current public vertex and $Z$. Its stored program can contain arbitrary read-only model constants but cannot recover the realized branch from those constants. An inherited intermediate tier, a customer identifier with a branch lookup, a retained shared random seed, and a transmitted numerical payment are all additional history channels; none is free in this interface. An $m$-symbol alphabet requires $\lceil\log_2m\rceil$ fixed-width writable bits, not $m$ bits. Public-state storage and read-only program storage are counted separately.

Private encoder randomization is allowed under expected participation. Fresh decoder randomness is independent of earlier private choices conditional on the symbol. Any seed that correlates those choices across dates must itself be encoded in the charged symbol. Concavity permits a deterministic terminal decoder at the expected-participation optimum, so a codebook is a collection of terminal levels. The lower bound for exact implementation continues to allow randomized encoders and decoders under this accounting. Distinct full-information terminal actions and strict concavity remain visible assumptions; Proposition~\ref{prop:repeated} handles repeated caps; uniqueness is required for the exact-action lower bound.


\subsection{Acceptance timing}
The renewal desk observes the branch and offers the intermediate tier and terminal lottery. Under expected participation, the customer accepts before the ticket draw; after acceptance only the ticket reaches the common gateway. The realized ticket may be disclosed, but no new costless exit option is supplied after the draw within this episode. Under pathwise participation, each possible realized obligation must instead satisfy the branch cap, and a draw-specific intermediate tier is allowed. A seed known before acceptance changes that conditioning information. Repeated-customer effects require an enlarged state and are not part of the three-date theorem. The full stipulated mechanism and its qualifications are given in the electronic companion; no empirical or legal claim about an actual provider is made. Section~\ref{sec:promise} allows a nonsaturated root promise, and Section~\ref{sec:catalog} specifies the intermediate risk protocol.
